\chapter{場の高速化と二軸周期境界}
\label{ch:field}

\section{Direct、treecode、FMM の役割}
$N$ source に対する $M$ 点の Direct 評価は $O(MN)$ の panel 和である。
同じ P0 source のまま遠方近似を導入しないため、小規模の基準計算として利用できる。
全要素重心の電位履歴を Direct で評価すると $O(N^2)$ になり、
粒子追跡とは別に出力が計算時間を支配する可能性がある。

treecode は source の octree を評価点ごとに走査し、十分に遠い node を一つの monopole で置き換える。
FMM は target 側にも局所展開を作り、遠方群の寄与を多くの評価点で共有する。
これらの手法の基礎は Barnes--Hut および Greengard--Rokhlin にある\citep{BarnesHut1986,GreengardRokhlin1987}。
BEACH の \code{auto} は要素数から Direct または FMM を選ぶ。
treecode の通常の対応場は自由空間であり、周期場では Direct/FMM と周期 split の互換性を確認する。

\section{treecode の遠方判定と電荷相殺}
node 内の総電荷、絶対電荷和、電荷中心を
\begin{equation}
 Q_n=\sum_{i\in n}q_i,\quad S_n=\sum_{i\in n}|q_i|,\quad
 \vect c_{Q,n}=\frac{\sum_{i\in n}q_i\vect c_i}{Q_n}
\end{equation}
とする。$Q_n$ が極小のときは幾何中心を用いる。
node 半径 $R_n$ は全三角形頂点を覆い、評価点との距離 $d_n$ に対して
\begin{equation}
 R_n<\theta(d_n-R_n)
\end{equation}
を遠方候補の条件とする。重心だけを用いて半径を過小評価すると、
評価点に近い大きな panel を誤って monopole にまとめる可能性がある。

さらに、正負電荷の相殺による電荷中心の不安定化を避けるため、
\begin{equation}
 |S_n-|Q_n||\le64\epsilon_{\mathrm{mach}}\max(S_n,|Q_n|)
\end{equation}
を満たす node だけを受理する。実質的に異符号電荷が混在する node では子へ降りる。
近傍と leaf は解析 panel 和で評価する。
この規則は相殺時の精度を守る一方、細かく正負電荷が混在する分布では高速化効果を弱める。

\section{Cartesian FMM の展開}
多重指数 $\alpha=(\alpha_x,\alpha_y,\alpha_z)$ を用い、
$|\alpha|=\alpha_x+\alpha_y+\alpha_z$、
$\alpha!=\alpha_x!\alpha_y!\alpha_z!$ とする。
中心 $\vect c$ の panel moment を、階乗を含む規約で
\begin{equation}
 M_\alpha(\vect c)
 =\sum_i\frac{q_i}{A_i\alpha!}
   \int_{T_i}(\vect y-\vect c)^\alpha\dd S_{\vect y}
 \label{eq:panel_moment}
\end{equation}
と定義する。重心点の monomial ではなく、三角形上の厳密な面積平均を用いる。
総電荷 $q_i$ には面積を再び掛けない。

FMM の手順は、leaf の moment を作る P2M、親へ集約する M2M、
遠方 source を target の局所展開へ移す M2L、子へ伝える L2L、
評価点で計算する L2P からなる。近傍は解析 panel の Direct 和を加える。
例えば、中心間変位を $\vect d=\vect c_{\mathrm{child}}-\vect c_{\mathrm{parent}}$ とすると
\begin{equation}
 M_\beta(\vect c_{\mathrm{parent}})
 =\sum_{\mathrm{child}}\sum_{\alpha\le\beta}
 M_\alpha(\vect c_{\mathrm{child}})
 \frac{\vect d^{\beta-\alpha}}{(\beta-\alpha)!}
\end{equation}
で集約する。$G(\vect r)=1/\lVert\vect r\rVert$、
$\vect R=\vect c_t-\vect c_s$ に対して M2L は
\begin{equation}
 L_\alpha(\vect c_t)\mathrel{+}=
 \sum_\beta(-1)^{|\beta|}M_\beta(\vect c_s)
 \partial^{\alpha+\beta}G(\vect R)
\end{equation}
となる。$L_\alpha$ は局所電位の微分係数であり、
\begin{equation}
 \phi_{\mathrm{far}}(\vect x)\simeq
 \kc\sum_{|\alpha|\le p}L_\alpha
 \frac{(\vect x-\vect c_t)^\alpha}{\alpha!},\qquad
 \vect E_{\mathrm{far}}=-\nabla\phi_{\mathrm{far}}
\end{equation}
で評価する。

固定 geometry の木、相互作用 pair、平行移動係数を初期化時に作り、
電荷に依存する moment と local だけをバッチごとに更新する。
現行 simulator の展開次数は $p=4$ 固定である。
低水準 FMM test が次数を変えられることと、\code{beach.toml} に公開次数キーがあることは異なる。
研究ケースでは \code{tree_theta}、leaf size、Direct 基準との差を確認する。
一定次数で良好な木分割を仮定した FMM の計算量の議論を、任意のケースの実測速度として引用しない。

\section{周期場の非零モードと零モード}
\code{periodic2} は $x,y$ 周期、$z$ 非周期の slab を対象とする。
mesh は primitive cell の base triangle を保持し、場と衝突は必要な画像を扱う。
粒子の周期 wrap を指定することと、静電場を無限周期にすることは別である。

cell の横面積を $A_{\mathrm c}=L_xL_y$ とすると、場を
\begin{equation}
 K_{\mathrm s}=
 \underbrace{K_{\mathrm{shell}}+R_{\mathrm{Ewald}}^{\mathrm{full}}
                 -K_0^{\mathrm{sym}}}_{k\ne0\ \text{の backend}}
 +K_0^{\mathrm{physical}}
 \label{eq:periodic_split}
\end{equation}
へ分ける。$K$ は電位と電場の評価作用素を表す。
有限画像だけの設定では、指定した shell の外側を含まない。
\code{field_periodic_far_correction="none"}、および現行の \code{auto} は有限画像モデルである。
無限周期 nonzero mode の production 経路は \code{cached_kneq0} を明示する。
\code{panel_spectral_reference} は小規模な P0 参照経路として用いる。

\section{Ewald 参照から生成する遠方 operator}
Ewald 分割は
\begin{equation}
 \frac1r=\frac{\operatorname{erfc}(a r)}r
            +\frac{\operatorname{erf}(a r)}r
\end{equation}
によって実空間と逆空間の収束を分担する\citep{LindboTornberg2012}。
$a$ は数値計算上の分割 parameter であり、Debye 遮蔽長の逆数ではない。
BEACH の teacher は実・逆空間を有限層で打ち切った Ewald2P 参照である。

初回生成では source の proxy 点と target の check 点を置き、
teacher から有限画像分を引いた残差を再現する operator を正則化 QR で求める。
root multipole を $\vect M_{\mathrm{root}}$ とすると、target anchor $t$ の補正は
\begin{equation}
 \vect L_t^{\mathrm{far}}=\mathsf A_t\vect M_{\mathrm{root}}
\end{equation}
である。電場 fit では決まらない電位の定数係数を、電位残差から別に fit する。
生成した operator を geometry・設定の fingerprint と checksum 付きで保存し、後の run で再利用する。
通常のバッチでは ordinary M2L の後、L2L の前に補正を加えるため、粒子ごとに Ewald 和を再計算しない。

確認した作業ツリーには、高次 moment 行の長さ尺度を揃える生成器修正が含まれる。
それでも、大きな有限画像場と零モードを差し引いて小さな非零場を得る領域では、残差誤差が支配し得る。
cache の再利用と単位尺度の一致だけで、弱い外部場に対する絶対精度を保証できない。

\section{Gauss 則で閉じる物理的零モード}
三角形 $i$ のうち高さ $z$ 以下の面積割合を $F_i(z)$ とおくと、累積電荷は
\begin{equation}
 C(z)=\sum_iq_iF_i(z)
\end{equation}
である。傾いた三角形の $F_i$ は頂点高さの間で区分二次、水平三角形は sheet の跳びになる。
BEACH はこの幾何係数を前計算し、電荷変更時に係数と積分 primitive を更新する。
下側遠方電場 $E_{\mathrm b}$ と gauge $(z_g,\phi_g)$ を指定すると
\begin{align}
 E_{0,z}(z)&=E_{\mathrm b}+\frac{C(z)}{\eps A_{\mathrm c}},
   \label{eq:zero_field}\\
 \phi_0(z)&=\phi_g-E_{\mathrm b}(z-z_g)
   -\frac1{\eps A_{\mathrm c}}\int_{z_g}^{z}C(\zeta)\dd\zeta
\end{align}
となる。全電荷 $Q=\sum_iq_i$ に対する二つの closure は
\begin{equation}
 E_{\mathrm b}=
 \begin{cases}
 -Q/(2\eps A_{\mathrm c}) & \text{対称な上下真空},\\
 0 & \text{下側電束を零に指定}
 \end{cases}
\end{equation}
である。どちらも材料内部の誘電分極や遮蔽を解くモデルではない。
電荷源より上では $C(z)=Q$ を用い、多項式の丸め残差を上部真空へ外挿しない。

非中性 cell の遠方には一定電場と線形電位が残り得る。
また $Q=0$ でも、高さによって正負の電荷が分離すれば $C(z)$ は零とは限らない。
零モードを単に消すと、この高さ分布と Gauss 則を失う。
外部シースを接続する場合も、表面電荷の零モードと外部空間電荷の役割を分け、二重加算を避ける。

\basisdoc{\src{docs/Treecode.md}、\src{docs/FMM.md}、\src{docs/FMMCore.md}、
\src{docs/PeriodicElectrostatics.md}、\src{docs/PeriodicFarCorrection.md}。}
